%%%%%%%%%%%%%%%%%%%%%%%%%%%%%%%%%%%%%%%%%
% "ModernCV" CV and Cover Letter
% LaTeX Template
% Version 1.1 (9/12/12)
%
% This template has been downloaded from:
% http://www.LaTeXTemplates.com
%
% Original author:
% Xavier Danaux (xdanaux@gmail.com)
%
% License:
% CC BY-NC-SA 3.0 (http://creativecommons.org/licenses/by-nc-sa/3.0/)
%
% Important note:
% This template requires the moderncv.cls and .sty files to be in the same 
% directory as this .tex file. These files provide the resume style and themes 
% used for structuring the document.
%
%%%%%%%%%%%%%%%%%%%%%%%%%%%%%%%%%%%%%%%%%

%----------------------------------------------------------------------------------------
%	PACKAGES AND OTHER DOCUMENT CONFIGURATIONS
%----------------------------------------------------------------------------------------

\documentclass[11pt,a4paper,sans]{moderncv} % Font sizes: 10, 11, or 12; paper sizes: a4paper, letterpaper, a5paper, legalpaper, executivepaper or landscape; font families: sans or roman
\usepackage{standalone}
% Load color BEFORE style so derived colors pick up blue values
\moderncvcolor{blue}
\moderncvstyle{casual}

\renewcommand*{\labelitemi}{\strut{\tiny\textbullet}}

% in the preamble
\usepackage{xcolor}
\definecolor{cvblue}{HTML}{2B6EA3} % replace with your CV hex
\usepackage{hyperref}
\hypersetup{
  colorlinks=true,
  linkcolor=cvblue,
  urlcolor=cvblue,
  citecolor=cvblue,
  pdfborder={0 0 0} % remove boxed borders
}

\colorlet{cvsubtle}{cvblue!60!black}
\hypersetup{urlcolor=cvsubtle, linkcolor=cvsubtle, citecolor=cvsubtle}


\usepackage[scale=0.85]{geometry} % Reduce document margins
%\setlength{\hintscolumnwidth}{3cm} % Uncomment to change the width of the dates column



%----------------------------------------------------------------------------------------
%	NAME AND CONTACT INFORMATION SECTION
%----------------------------------------------------------------------------------------

\name{Edouard}{Pignède}

% All information in this block is optional, comment out any lines you don't need
\title{Curriculum Vitae}
% \address{Paris-Saclay Applied Economic (PSAE)}{12 place de l'Agronomie, 91120 Palaiseau}
 %\social{github}{stefano-bragaglia}
\email{edouard.pignede@gmail.com} 



\homepage{edouardpignede.github.io}

\social[github]{edouardpignede}
\social[linkedin]{edouard-pignede}



%\social[linkedin][www.linkedin.com]{name}
% The first argument is %the url for the clickable link, the second argument is the url displayed in the %template - this allows special characters to be displayed such as the tilde in this %example

\photo[70pt][0.3pt]{picture} % The first bracket is the picture height, the second is %the thickness of the frame around the picture (0pt for no frame)
%\quote{Not Attention, Patience is all we need.}

%----------------------------------------------------------------------------------------

\newcommand{\cvdoublecolumn}[2]{%
  \cvitem[.75em]{}{%
    \begin{minipage}[t]{\listdoubleitemcolumnwidth}#1\end{minipage}%
    \hfill%
    \begin{minipage}[t]{\listdoubleitemcolumnwidth}#2\end{minipage}%
    }%
}



\usepackage{multibbl}
\newcommand\Colorhref[3][orange]{\href{#2}{\small\color{#1}#3}}


% \newcommand{\cvreference}[7]{%
%     \textbf{#1}\newline% Name
%     \ifthenelse{\equal{#2}{}}{}{\addresssymbol~#2\newline}%
%     \ifthenelse{\equal{#3}{}}{}{#3\newline}%
%     \ifthenelse{\equal{#4}{}}{}{#4\newline}%
%     \ifthenelse{\equal{#5}{}}{}{#5\newline}%
%     \ifthenelse{\equal{#6}{}}{}{\emailsymbol~\texttt{#6}\newline}%
%     \ifthenelse{\equal{#7}{}}{}{\phonesymbol~#7}}

\begin{document}

\makecvtitle % Print the CV title




%----------------------------------------------------------------------------------------
%	EDUCATION SECTION
%----------------------------------------------------------------------------------------
\section{Current position}

\cventry{2025--Present}{Post-doctoral scholar}{\href{https://en.ird.fr/}{French National Research Institute for Sustainable Development (IRD), \href{https://www.umi-source.uvsq.fr}{Université Paris-Saclay, UMI SOURCE}}}{Paris}{}
{
Member of the BETSAKA project, initiative to assess the environmental and socio-economic impacts of protected areas in Madagascar from 2000 to 2024. \newline
3-month fieldwork in Madagascar (March to June 2026) to organize interviews with protected area managers and stakeholders.
}

\section{Academic appointments}

\cventry{2025}{Research Fellow}{\href{https://novafrica.org/}{NovAfrica} (\href{https://www.novasbe.unl.pt/en/}{NOVA SBE})}{Lisbon}
{\newline Member of the FCTPex project \textit{Artisanal Gold Mining and Sustainable Development}}{}

\section{Education}

\cventry{2021--2025}{PhD, Development and environmental economics}{\href{http://www2.agroparistech.fr/spip.php?article1142}{AgroParistech} (\href{https://www.universite-paris-saclay.fr/en}{Paris-Saclay University}) \& \href{https://www.chaireeconomieduclimat.org/en/}{Climate Economic Chair}}{Paris}{}
{
    \href{https://www.afse.fr/en/news/winners-of-the-2026-afse-thesis-prize-2531}{2026 Thesis Award of the French Economic Association (AFSE)} \newline 
    Title: The poverty-environment trap : Climate change, natural resources and population dynamics in Sub-Saharan Africa (\href{https://theses.fr/2025UPASB038}{link}) \newline
    Supervisors: \href{https://sites.google.com/view/philippedelacote}{Philippe Delacote}, \href{https://sites.google.com/site/rajachakir/}{Raja Chakir} and \href{https://julienwolfersberger.fr/}{Julien Wolfersberger}
    \begin{itemize}
        \item[] Chapter 1: impact of drought on inequalities using quantile treatment effect and distribution regression, as well as the initial endowment in capital.
        \item[] Chapter 2: empirical analysis of the link between artisanal mining and urbanization.
        \item[] Chapter 3: quantitative spatial modeling of climate change-induced migration with liquidity constraints.
        \item[] Chapter 4: quantitative spatial modeling of climate change impact on mortality.
    \end{itemize}
}  % Arguments not required can be left empty

\cventry{2019--2020}{Master’s degree in environmental economics}{\href{http://www2.agroparistech.fr/spip.php?article1142}{AgroParistech} (\href{https://www.universite-paris-saclay.fr/en}{Paris-Saclay University}) }{Paris}{}
{
    Average: 15.95 \newline
    Courses: environmental econometrics, economy of natural resources, international trade and environment, public policies and environment, climate change economy
}
%{Advanced exposure to various areas of computer science along with a one and half year research project on Reversible Logic Synthesis.}
%\cvitem{CGPA :}{7.96/10}

\cventry{2016--2020}{Master’s degree in industrial engineering}{\href{https://www.centralesupelec.fr/en}{CentraleSupélec} (\href{https://www.universite-paris-saclay.fr/en}{Paris-Saclay University})}{Paris}
{}
{
    Average: 16.35
    \begin{itemize}
        \item[] Specialization in industrial engineering and management of international project.
        \item[] Courses: advanced probability, advanced statistics, economy, sustainable development, modelling and optimization.
        \item[] Six-month project on applied mathematics: development of modelling algorithm of Hawkes process computed with Python.
    \end{itemize}
}


\cventry{2014--2016}{Classes préparatoires}{Henri IV}{Paris}{}
{A two-year intensive program in mathematics and physics to prepare for nationwide competitive entrance in top French graduate schools of engineering.}


%{Comprehensive exposure to the core areas of Computer Science along with a final year project on Data-mining}
%\cvitem{CGPA :}{7.36/10}
% \cventry{2008 :}{Higher Secondary Examination}{Belmuri Union Institution}{Belmuri}{}{ Mathematics, Physics, Chemistry, Biology, English, Bengali}
% {}
% \cvitem{Percentage :}{81.2 \%}
% \cventry{2006 :}{Secondary Examination}{Belmuri Union Institution}{Belmuri}{}{ Mathematics, Physical Science, Life Science, Geography, History, English, Bengali}
% {}
% \cvitem{Percentage :}{90.8 \%}





\newpage
%----------------------------------------------------------------------------------------
%	PUBLICATION SECTION
%----------------------------------------------------------------------------------------


\section{Publications}
% \subsection{Journal Article(Accepted)}
% \cventry{2019}{\textbf{Pratik Dutta}, Sriparna Saha, Sanket Pai and Aviral Kumar}{}{Protein-protein Interaction based Generative Model for Improving Gene Clustering}{In \textit{\textbf{Scientific Reports-Nature}} (\textbf{Impact Factor: 4.12)}}{}

\subsection{Journal Articles}
\newbibliography{journal}
\bibliographystyle{journal}{plainyrrev}
\nocite{journal}{*}
\bibliography{journal}{journal}
{\large \textsc{Refereed Journal Articles}}

\subsection{Non-refereed journal articles}
\newbibliography{unpublished}
\nocite{unpublished}{*}
\bibliographystyle{unpublished}{plainyrrev}
\bibliography{unpublished}{unpublished}
{\large \textsc{Refereed Conference Publications}}

\subsection{On-going works}
\cventry{2026}{Edouard Pignède}{}{Florent Bedecarrats and Thomas Thivillon, Protected areas and fires in Madagascar}{}{}
\cventry{2025}{Edouard Pignède}{}{Climate change, adaptation, and mortality}{}{}
\cventry{2024}{Edouard Pignède}{}{and Julien Wolfersberger, Climate immobility in sub-Saharan Africa}{Available uppon request}{}

%\subsection{Communicated Journal Article}
%\cventry{2020}{Pratik Dutta, Aditya Prakash Patra, and Sriparna Saha}{}{DeePROG: An Attention based Deep Multi-modal Architecture for Disease Gene Prognosis}{In \textit{IEEE Transactions on Biomedical Engineering}}{}


% \subsection{Conference and talks}
% \newbibliography{conference}
% \nocite{conference}{*}
% \bibliographystyle{conference}{plainyrrev}
% \bibliography{conference}{conference}
% {\large \textsc{Refereed Conference Publications}}









%----------------------------------------------------------------------------------------
%	WORK EXPERIENCE SECTION
%----------------------------------------------------------------------------------------

\section{Research Experience}
\cventry{Nov,2020 -- Nov,2021}{\textit{Research engineer}}{\href{https://en.ird.fr/}{French National Research Institute for Sustainable Development (IRD), \href{https://dial.ird.fr/en/}{DIAL}}}{Paris}{}
{
    \begin{itemize}
        \item[] Mapping of Sub-Saharan vulnerability to climate change through three components: exposition, sensibility and adaptation capacities. Building of multiple indicators of vulnerability using household survey (DHS, agricultural survey), satellite data (CHIRPS, ISRIC) and country data (infrastructures, mines localisation) computed with RStudio and QGIS.
        \item[] Application of an original statistical method of spatial matching to assess the effect of roads on inter-sectoral mobility and population migration in Mali, using national census data and RStudio.
        \item[] Building a global database of coast escarpment, nightlight and wave height to assess the role of piracy on coast development using RStudio and QGIS.
        \item[] Computing of price indexes in France in the 18th century using the repeated sales method, RStudio and data from the TOFLIT project (published in Dialogue n°62, English version: \href{https://dial.ird.fr/wp-content/uploads/2021/10/Dialogue_Eng_62_def.pdf}{link}. French version: \href{https://dial.ird.fr/wp-content/uploads/2021/10/Dialogue-62.pdf}{link}).
    \end{itemize} 
}

\cventry{May,2020 -- Nov,2020}{\textit{Research engineer}}{\href{https://www.afd.fr/en}{French Development Agency (AFD)}}{Paris}{End-of-study internship}
{\begin{itemize}
    \item[] Development of statistical models of prediction of sugarcane yields in Ivory Coast based on machine learning algorithm (random forest) and using meteorological data (precipitations, temperatures) and satellite data (NDVI, EVI). Computation with RStudio and QGIS.
    \item[] Publication of a scientific paper in a peer-reviewed journal as well as a research paper in French.
\end{itemize} }


%\cvitem{Advisor :}{\textbf{Dr. abc xyz}, \textit{Associate Professor, Department of Computer Science \& Engineering}, IIT abc ({\Colorhref{https://www.personal_webpage.com/} {\textit{Personal Web-page}}})}


\section{Conference and talks}

\cvitem{2026}{
    CLAND Workshop - Land Use Policies for Climate Change, Paris (France) \newline
    ADRES - Job Market Conference, Paris (France) \newline
    Seminar of UMI SOURCE, University of Paris-Saclay
}

\cvitem{2025}{
    FAERE - Conference of the French Association of Environmental Economists, Nantes (France).\newline
    EEA Congress, Bordeaux (France). \newline
    ICDE - International Conference in Development Economics, Nanterre (France).\newline
    CEEM Seminar, Agro Institute of Montpellier - INRAE.\newline
    CSAE Conference in Development Economics, Oxford (England).\newline
    Lisbon Micro Group, ISEG, Lisbon (Portugal).\newline
    DIAL Seminar, Paris-Dauphine University.
}

\cvitem{2024}{
    EAERE - Conference of the European Association of Environmental Economics, Leuven (Belgium).\newline
    NOVAFRICA Conference on Economic Development, Lisbon (Portugal).\newline
    Paris-Saclay Conference on Trade and the Environment, University Paris-Saclay \newline
    CSAE Conference in Development Economics, Oxford (England).\newline
    Workshop in Environmental Economics, Toulouse School of Economics (France).
}

\cvitem{2023}{
    PhD Research Group of NOVAFRICA, Lisbon (Portugal).\newline
     Lisbon Micro Group, ISEG, Lisbon (Portugal).\newline
     FAERE - Conference of the French Association of Environmental Economists, Montpellier (France).\newline
     EAAE - European Association of Agricultural Economists Congress, Rennes (France).\newline
    ICDE - International Conference in Development Economics, Paris (France).\newline
    AFSE - Congress of the French Economic Association, Paris (France).\newline
    Paris-Saclay Applied Economics Internal Seminar, Paris (France).\newline
    EAERE Winter School on Natural Resources and Development, Ascona (Switzerland).}

\cvitem{2022}{
    AgEconMeet - Early career seminar for agricultural economists in Europe, Göttingen (Germany).
}



% \section{Other work experience}

% \cventry{Sep,2017 -- Sep,2019}{\textit{Treasurer and member of LNA  project}}{}{}{}
% {\begin{itemize}
%     \item Member of a solidarity project aiming at offering musical workshops to disadvantaged children during six months in Burkina Faso, Vietnam, Cambodia and Bolivia.
%     \item In charge of a 50 000€ budget, workshop teaching, music composition and concert management.
% \end{itemize}}

% \cventry{Jul,2018 -- Dec,2018}{\textit{Data analyst}}{\href{https://alphalyr.fr/}{Alphalyr}}{Paris}{}
% {\begin{itemize}
%     \item Evaluation of advertising campaigns from the data obtained by the SaaS software Elitrack.
%     \item Creation of an automatic tool to achieve this evaluation.
% \end{itemize}}













%----------------------------------------------------------------------------------------
%	Fellowships \& Awards
%----------------------------------------------------------------------------------------

%\section{Fellowships \& Awards}

%\cvitem{2016 --present}{\textit{\textbf{Visvesvaraya Fellowship}} of Ministry of Electronics and Information Technology (MeitY), Government of India, as a PhD research scholar in Indian Institute of Technology Patna.}
%\cvitem{2019}{Receipt of \textit{\textbf{Visvesvaraya Travel Grant}} to attend a international conference \textbf{\textit{IEEE Congress on Evolutionary Computation, 2019}} in Wellington, New Zealand.}
%\cvitem{2018}{Recipient of \textit{\textbf{SciGenome Research Foundation (SGRF) GYAN Scholarship}} to participate \textbf{\textit{Nextgen Genomics, Biology, Bioinformatics and Technologies-2018}} meeting at Jaipur India from $30^{th}$ September to $2^{nd}$ October 2018. }
%\cvitem{2015}{Awarded under \textit{\textbf{Students Reward Programme}} at the Annual General Meeting of \textbf{Global Alumni Association of Bengal Engineering and Science University(GAABESU).}}


%----------------------------------------------------------------------------------------
%	Academic achievements
%----------------------------------------------------------------------------------------

%\section{Academic Achievements \& Recognitions }


%\cvitem{2018}{\textbf{Session Chair of the session "Prediction"} in \textbf{$25^{th}$ \textit{International Conference of Neural Information Processing} (ICONIP 2018)}, Siem Reap, Cambodia.}


%\cvitem{2018}{Invited to conduct lab sessions in \textit{\textbf{"Training Program on Machine Learning For Ocean Acoustics and Climate Data Analysis"}}, during 22-36 October 2018 at \textbf{Defence R\&D Organization- Naval Physical \& Oceanographic Laboratory (DRDO-NPOL), Kochi, Kerala}.}


%----------------------------------------------------------------------------------------
%	COMPUTER SKILLS SECTION
%----------------------------------------------------------------------------------------

\section{Skills}

\cvitem{Computer skills}{Advanced knowledge of R, Python and QGIS
\newline Environment: Quarto, Positron, SSP Cloud, GitHub, Zotero
\newline Basic knowledge of STATA}

\cvitem{Languages}{French: Native
\newline English: Fluent
\newline Spanish: Intermediate}


%----------------------------------------------------------------------------------------
%	Position of Responsibility SECTION
%----------------------------------------------------------------------------------------

%\section{Position of Responsibility}
%\cventry{2016-2020}{Executive member of IEEE Student Branch}{}{IIT ABC}{}{}
%\cventry{April 1-5, 2019}{Organizer, GIAN Workshop on subjects}{}{IIT ABC}{}{}




%----------------------------------------------------------------------------------------
%	Teaching  SECTION
%----------------------------------------------------------------------------------------

\section{Teaching experience}

\cventry{2026}{Training in research replicability}{PhD Program - \href{https://www.umi-source.uvsq.fr/cv-doctorants-7}{CERED} (\href{https://www.univ-antananarivo.mg/}{University of Antananarivo})}{Madagascar}{}{Taught in French. 14-hour course at the PhD level.}

\cventry{2023}{Training in research replicability}{PhD Program - \href{http://www2.agroparistech.fr/spip.php?article1142}{PSAE} (\href{https://psae-saclay.fr/}{Paris-Saclay University})}{Paris}{}{Taught in French. 2-hour course at the PhD level.}

\cventry{2023}{Introduction to Environmental Economics}{Master BIOCEB - \href{http://www2.agroparistech.fr/spip.php?article1142}{AgroParistech} (\href{https://www.universite-paris-saclay.fr/en}{Paris-Saclay University})}{Paris}{}{Taught in English. 9-hour course at the Master level.}

\cventry{2022}{Introduction to Climate Change in Economics}{Master ADEPP - \href{https://www.ensae.sn/accueil}{ENSAE Daka}r}{Dakar}{}{Taught in French. 2-hour course at the Master level.}

\cventry{2022}{Microeconomy}{CPES - \href{https://psl.eu/en}{PSL University}}{Paris}{}{Taught in French. 20-hour course at the Bachelor level (2nd year).}


\cventry{2022}{On the use of GIS data in economy}{\href{http://www2.agroparistech.fr/spip.php?article1142}{AgroParistech} (\href{https://www.universite-paris-saclay.fr/en}{Paris-Saclay University})}{Paris}{}{One session introduction lecture course to M2 students.}
\cventry{2018}{Teaching assistant in Information System}{\href{https://www.centralesupelec.fr/en}{CentraleSupélec} (\href{https://www.universite-paris-saclay.fr/en}{Paris-Saclay University})}{Paris}{}{Course for 1st year student.}
\cventry{2018}{Personal mathematics teacher}{}{Paris}{}{Teaching mathematics to 10 high-school students.}



\newpage
%----------------------------------------------------------------------------------------
%	Referee reports
%----------------------------------------------------------------------------------------

\section{Referee reports for scientific journals}

\cvitem{2026}{IPCC Special Report on Climate Change and Cities}{}{}
\cvitem{2025}{American Journal of Agricultural Economics}{}{}
\cvitem{}{Agricultural Economics}{}{}
\cvitem{}{Journal of African Economies}{}{}
\cvitem{2024}{American Journal of Agricultural Economics}{}{}



%----------------------------------------------------------------------------------------
%	References
%----------------------------------------------------------------------------------------


\section{References}


\begin{tabular}{lr}
% Referee 1
\begin{minipage}[t]{3.5in}
\textbf{\href{https://sites.google.com/view/philippedelacote}{Dr. Philippe Delacote}}\\
\textit{Senior Researcher, Department of Economy} \\
\textit{Bureau d'Economie Théorique et Appliquée (BETA)}\\
INRAE \\
\faEnvelope\ \href{mailto:philippe.delacote@nancy.inra.fr}{philippe.delacote@inrae.fr}
\end{minipage}
&
% Referee 2
\begin{minipage}[t]{3.5in}
\textbf{\href{https://sites.google.com/site/rajachakir/}{Dr. Raja Chakir}}\\
\textit{Senior Researcher, Department of Economy} \\
\textit{Paris-Saclay Applied Economics (PSAE)}\\
INRAE-AgroParisTech\\
\faEnvelope\ \href{mailto:raja.chakir@inrae.fr}{raja.chakir@inrae.fr}
\end{minipage}
\\
\\ % Additional newline for spacing.
% Referee 3
\begin{minipage}[t]{3.5in}
\textbf{\href{https://julienwolfersberger.fr/}{Dr. Julien Wolfersberger}}\\
\textit{Research fellow, Department of Economy} \\
\textit{Paris-Saclay Applied Economics (PSAE)}\\
AgroParisTech\\
\faEnvelope\ \href{mailto:julien.wolfersberger@agroparistech.fr}{julien.wolfersberger@agroparistech.fr}
\end{minipage}
&
\begin{minipage}[t]{3.5in}
\textbf{\href{https://sites.google.com/site/girardvictoire/home}{Dr. Victoire Girard}}\\
\textit{Research faculty, Department of Economy} \\
\textit{NovAfrica}\\
 Nova SBE\\
\faEnvelope\ \href{mailto:victoire.girard@novasbe.pt}{victoire.girard@novasbe.pt}
\end{minipage}
\\
\end{tabular}


\end{document}